\documentclass{article}
\usepackage[utf8]{inputenc}
\usepackage{amsmath}
\usepackage{geometry}
\geometry{margin=2cm}
\begin{document}

\section*{Análise Completa da Questão 147 — ENEM 2024 — Matemática}

\section{Solução Técnica}
A alternativa correta é \textbf{B) II, IV e V}.

\textbf{Justificativa:} Os números nas posições II, IV e V representam a mesma quantidade.

\textbf{Estratégia:} Converter todos os números apresentados nas posições I a VI para uma forma comum, preferencialmente decimal, e comparar seus valores para identificar quais são iguais.

\textbf{Passos:}
\begin{itemize}
  \item Converter os números nas posições para forma decimal:
  \begin{align*}
    \text{I: } & 4 \frac{1}{2} = 4 + \frac{1}{2} = 4,5 \
    \text{II: } & 4 \frac{1}{2} = 4 + \frac{1}{2} = 4,5 \
    \text{III: } & \frac{10}{45} = \frac{2}{9} \approx 0,222 \
    \text{IV: } & \frac{18}{4} = 4,5 \
    \text{V: } & 4,5 \quad (\text{já decimal}) \
    \text{VI: } & \frac{4}{5} = 0,8
  \end{align*}
  \item Comparar os valores obtidos. As posições I, II, IV e V têm valor igual a 4,5, enquanto III e VI representam valores diferentes.
  \item A questão solicita três quantidades iguais; o estudante escolheu as posições II, IV e V, que representam 4,5.
\end{itemize}

\section{Fórmulas Utilizadas}
\begin{itemize}
  \item Conversão de número misto para decimal:
    \[
      a \; \frac{b}{c} = a + \frac{b}{c}
    \]
  \item Simplificação de frações:
    \[
      \frac{10}{45} = \frac{2}{9}
    \]
  \item Comparação de números decimais para identificar valores equivalentes.
\end{itemize}

\section{Evidências Visuais Utilizadas}
A questão apresenta uma imagem de uma máquina com um visor contendo seis compartimentos numerados de I a VI, contendo os seguintes números:
\begin{itemize}
  \item I: $4 \frac{1}{2}$
  \item II: $4 \frac{1}{2}$
  \item III: $\frac{10}{45}$
  \item IV: $\frac{18}{4}$
  \item V: $4,5$
  \item VI: $\frac{4}{5}$
\end{itemize}
Essas posições numéricas e seus valores foram essenciais para realizar a conversão e comparação para solução correta da questão.

\section{Explicação Didática}

\subsection{Explicação Resumida}
Para determinar quais números representam a mesma quantidade, convertemos todos para a mesma forma — no caso, números decimais — e comparamos seus valores. As posições II, IV e V têm valor 4,5, ou seja, representam a mesma quantidade.

\subsection{Passo a passo}
\begin{enumerate}
  \item Observe os números nas posições de I a VI.
  \item Converta cada número para decimal para facilitar a comparação:
  \begin{itemize}
    \item I: $4 \frac{1}{2} = 4 + 0,5 = 4,5$
    \item II: $4 \frac{1}{2} = 4,5$
    \item III: $\frac{10}{45} = \frac{2}{9} \approx 0,222$
    \item IV: $\frac{18}{4} = 4,5$
    \item V: $4,5$ (já decimal)
    \item VI: $\frac{4}{5} = 0,8$
  \end{itemize}
  \item Compare os valores decimais para encontrar os que são iguais.
  \item Identifique as posições II, IV e V como representando a mesma quantidade.
  \item Escolha essas posições como resposta à pergunta do exercício.
\end{enumerate}

\subsection{Erros comuns}
\begin{itemize}
  \item Não converter corretamente frações para decimais, causando erro na comparação.
  \item Escolher números que pareçam próximos, sem calcular o valor exato.
  \item Desconsiderar alguma posição que representa uma quantidade igual por falta de atenção na conversão.
  \item Pensar que números mistos com escrita diferente não podem ser iguais, sem verificar os valores decimais.
\end{itemize}

\subsection{Dicas para o ensino}
\begin{itemize}
  \item Incentive alunos a sempre converter diferentes representações numéricas para uma forma comum antes de comparar.
  \item Pratique conversões com números mistos, frações e decimais para fortalecer a compreensão.
  \item Ensine a simplificar frações para facilitar a comparação.
  \item Use recursos visuais ou manipulativos para mostrar equivalência entre diferentes representações.
  \item Oriente a investigar cuidadosamente os valores numéricos antes de tomar decisão na resposta.
\end{itemize}

\section{Erros comuns ("Common Pitfalls")}
\begin{itemize}
  \item Confundir ou não converter corretamente frações para decimais, levando a comparações incorretas.
  \item Selecionar números apenas por proximidade aparente, sem realizar cálculo exato ou conversão.
  \item Ignorar algum número que representa a mesma quantidade devido à falta de atenção nas conversões.
  \item Supor que números mistos com escrita diferente não possam ter valores iguais sem fazer a verificação decimal.
\end{itemize}

\section{Dicas para o Ensino}
\begin{itemize}
  \item Ensine os alunos a transformar diversas representações numéricas para uma forma padrão antes de comparar.
  \item Inclua exercícios de conversão entre números mistos, frações e decimais para reforçar o aprendizado.
  \item Explique a importância da simplificação de frações para facilitar identificações de equivalências.
  \item Utilize materiais visuais ou objetos concretos para exemplificar equivalências numéricas.
  \item Estimule os alunos a verificar cuidadosamente os valores numéricos e evitar respostas apressadas.
\end{itemize}

\section{Distratores e Análise}
\begin{itemize}
  \item \textbf{Alternativa A) I, II e IV}
  \begin{itemize}
    \item \textit{Raciocínio provável:} Pode parecer correto pois I, II e IV representam fracções e números mistos equivalentes a 4,5.
    \item \textit{Tipo de equívoco:} Confusão na quantidade solicitada ou seleção arbitrária sem considerar todas as opções.
    \item \textit{Por que parece plausível:} I, II e IV são números iguais a 4,5.
    \item \textit{Por que está errada:} A questão pede três quantidades iguais e a resposta correta é II, IV e V. Apesar de I ter o mesmo valor, o estudante não a escolheu.
  \end{itemize}
  \item \textbf{Alternativa C) II, III e V}
  \begin{itemize}
    \item \textit{Raciocínio provável:} Confundir o valor decimal da posição III, achando próximo de 4,5.
    \item \textit{Tipo de equívoco:} Não realizar a conversão correta de frações para decimal.
    \item \textit{Por que parece plausível:} III é uma fração, V decimal; podem parecer relacionados.
    \item \textit{Por que está errada:} III aprox. 0,222 é diferente de 4,5.
  \end{itemize}
  \item \textbf{Alternativa D) III, V e VI}
  \begin{itemize}
    \item \textit{Raciocínio provável:} Agrupamento incorreto sem análise adequada.
    \item \textit{Tipo de equívoco:} Erro na comparação numérica.
    \item \textit{Por que parece plausível:} Parece que os valores fracionários se relacionam, mas não são iguais.
    \item \textit{Por que está errada:} III $\approx 0,222$, V = 4,5, VI = 0,8 são valores distintos.
  \end{itemize}
  \item \textbf{Alternativa E) III, IV e VI}
  \begin{itemize}
    \item \textit{Raciocínio provável:} Confusão na equivalência das frações.
    \item \textit{Tipo de equívoco:} Erro de conversão e comparação.
    \item \textit{Por que parece plausível:} Pode parecer que 18/4 é próximo a 4/5, o que não é.
    \item \textit{Por que está errada:} IV = 4,5, VI = 0,8, III = 0,222 são diferentes.
  \end{itemize}
\end{itemize}

\section{Perfil da Questão}
\begin{itemize}
  \item \textbf{Habilidade principal:} Conversão e comparação de números racionais.
  \item \textbf{Habilidades secundárias:}
  \begin{itemize}
    \item Representação numérica.
    \item Simplificação de frações.
    \item Interpretação de informações visuais.
  \end{itemize}
  \item \textbf{Demanda cognitiva:} Média.
  \item \textbf{Dificuldade estimada:} Média.
  \item \textbf{Requer interpretação visual:} Sim.
  \item \textbf{Requer modelagem matemática:} Não.
  \item \textbf{Requer cálculos:} Sim.
  \item \textbf{Observações:}
  \begin{itemize}
    \item Necessário converter frações e números mistos para decimal para comparar corretamente.
    \item Atenção à representação de números racionais em diferentes formas.
    \item Uso da imagem para identificação dos números e suas posições.
  \end{itemize}
\end{itemize}


\end{document}